\section{Current scope: full-envelope identification}
\label{sec:current-scope}
\paragraph{Decision of 9 October 2026.}
The primary contribution is now full-envelope LPV identification from incomplete
local operating coverage. Measured speed represents operating variation;
learned compatible groups represent persistent client differences. The active
formulation assumes available states $x=[\beta,r]^\top$ for identification and
feedback. Scheduled LQI is supporting validation. The KF and higher-order
10H experiment chain is paused. This decision changes the manuscript scope;
it does not change any historical experiment or launch a new simulation.

\paragraph{Evidence to consolidate.}
Experiments 10A--10E provide the main chain: development-only representation
selection, complementary-coverage rank, finite noisy state-transition fitting,
learned compatibility and numerical federated equivalence. In the primary
label-free 10E setting, mean paired fleet full-envelope matrix-error reductions
are 93.41\% against restricted Local and 23.74\% against Global, with a
95\% fleet-bootstrap interval $[15.39,31.04]\%$ for the latter. The selected
models use two groups in eight of ten fleets and one in two. Group discovery
must therefore allow a shared global model as a valid outcome.

The 10F confirmation supports downstream model usability. Against Local,
tracking gains are 23.86\% for moderate and 8.96\% for hard maneuvers. Against
Global, the gains are 4.26\% and 1.21\%, respectively. The corresponding
absolute changes are only 0.001824 and 0.001874 degrees/s. A large tracking
advantage is not the primary
claim. The reference Global model here is fitted from fleet observations,
not the public nominal model used in reconstructed 10H-Q.

\paragraph{Assumptions and claim boundaries.}
The existing basis is the frozen orthogonalized span
$[1,v^{-1},v^{-2},z,z^2,z^3,z^4]$, $z=(v-20)/10$. Full envelope refers to
10--30 m/s along the audited cornering operating curve, not all nonlinear
state/input conditions. The finite-data benchmark uses independent state/input
perturbations around simulator-computed equilibria, exact current regressors
and noisy next-state targets. Group fitting receives estimated local matrices
and speed, but simulator-derived centering remains privileged setup information.
Consequently, this evidence does not yet validate learning from ordinary driving
records with no operating-point information.

Whole-envelope matrix errors must not be relabelled as unseen-speed-only errors.
The evaluated clients contribute restricted fitting records; fresh-fleet
replication is not frozen-model transfer to wholly held-out recipients.
Raw-data locality provides no formal privacy guarantee. Frozen-speed stability
checks and finite clipped trajectories provide no scheduled-stability certificate.
The active manuscript retains these limitations explicitly.

\paragraph{Disposition and next decision.}
The identification draft is \path{paper/main.tex}; the preceding cold-start
manuscript is preserved under \path{paper/archive/9m_cold_start/}. The detailed
claim/evidence map and one unresolved state-record interface boundary are in
the repository's consolidation note. No broad new campaign is planned.
If controlled equilibrium-centered data are retained, consolidate and review
the existing evidence. If a state-record-only learning claim is required, first
freeze one regression-interface gate using only supplied state/input/speed arrays
and empirically available centering, with the existing basis and identification
endpoint. Hidden physical parameters, true Jacobians and family labels must
remain outside that fit and selection.

The remainder of this report is the chronological development record. Older
scope statements and proposed next steps describe their original stage and
are superseded by this section for current work.
\clearpage
